\scanpage{79}
\mspage{\handwritten{24}}
D\handwritten{\'}emonstration. (i). Par r\handwritten{\'}ecurrence no\unsure{\handwritten{\'}e|e}th\handwritten{\'}erien\inserted{\handwritten{ne}} et (\handwritten{2.4 (i)}), il su\deleted{\unsure{ff}}ffit de d\handwritten{\'}emontrer que $f_{*}F'$ est localement consta\deleted{\unsure{tn}}nt et constructible sur un ouvert non-vi\deleted{\unsure{ed}}de \deleted{\handwritten{quelconque}}\inserted{\handwritten{convenable,}} C'est donc une assertion l\deleted{\unsure{i}}ocale sur $X$. Soit \deleted{\unsure{$x$}} $x$ un point maxi\deleted{\unsure{n}}mal de $X$. En rempla\handwritten{\c}cant $X$ par un rev\handwritten{\^}etement \handwritten{\'}etale d'un voisinage de $x$, on peut supposer que chaque point $x'$ de $X'$ au dess\deleted{\unsure{o}}us de $x$ est radiciel \deleted{\unsure{lui}}\handwritten{sur $x$}. \deleted{dessus\unsure{,|.}} Alors $s:X'\handwritten{\longrightarrow}X$ est somme de morp\deleted{g}hismes radiciels dans un voisinage de $x$ (\handwritten{\inserted{EGA} \uline{IV}.9.6.1}). \deleted{Rempla\handwritten{\c}cant} \inserted{\handwritten{Rétrécissant}} \deleted{\unsure{$X$}} encore\inserted{\handwritten{ $X$}}, on peut supposer que $f$ est somme \deleted{\unsure{e}}de morp\deleted{g}h\deleted{\unsure{e}}ismes radic\deleted{\unsure{leuses}}iels, et de plus que $F'$ est constant (et constructib\deleted{\unsure{e}}le) sur $X'$. Alors $f_{*}F'$ est produit d'un nombre fini de faisceaux constructibles constants, et est donc constructible\deleted{\unsure{s}} et constant.

\figure{\begin{tikzpicture}[x=1cm,y=1cm]
\draw[thin] (0,0) -- (-.04,.10) .. controls (1.80,.28) and (.57,-.42) .. (.23,-.90)
 .. controls (-.23,-1.58) and (1.80,-.97) .. (1.48,-1.91)
 .. controls (1.30,-2.48) and (.48,-2.72) .. (1.06,-3.14)
 .. controls (1.52,-3.43) and (1.95,-3.53) .. (1.84,-4.25)
 .. controls (1.70,-5.04) and (1.83,-5.19) .. (2.22,-5.72);
\end{tikzpicture}}
% archive sheet number: 78
